% 微分同胚（综述）
% license CCBYSA3
% type Wiki

本文根据 CC-BY-SA 协议转载翻译自维基百科\href{https://en.wikipedia.org/wiki/Diffeomorphism}{相关文章}

在数学中，微分同胚是可微流形之间的一种同构。它是一个可逆函数，能够将一个可微流形映射到另一个可微流形，并且该函数及其逆函数都是连续可微的。
\begin{figure}[ht]
\centering
\includegraphics[width=6cm]{./figures/61cddbbc19414b20.png}
\caption{一个从正方形映射到其自身的微分同胚下，正方形上矩形网格的映射图像。} \label{fig_WeiFTP_1}
\end{figure}
\subsection{定义}
给定两个可微流形 $M$ 和 $N$，一个连续可微映射$f: M \to N$如果它是双射，并且它的逆映射$f^{-1}: N \to M$也是可微的，那么 $f$ 就称为一个微分同胚。如果映射$f$及其逆映射都是$r$ 次连续可微的，则称 $f$是一个$C^r$-微分同胚。

两个流形$M$和$N$如果存在一个微分同胚$f: M \to N$则称它们是微分同胚的，通常记作$M \simeq N$。如果两个流形都是$C^r$可微的，并且存在一个$r$次连续可微的双射，其逆映射也是$r$次连续可微的，那么称它们是$C^r$-微分同胚的。$C^1$-微分同胚通常直接称为微分同胚。$C^0$-微分同胚则等价于一个同胚。
\subsection{流形子集的微分同胚}
设$M$和$N$是两个流形，$X \subset M$ 和 $Y \subset N$ 是它们的子集。一个函数$f: X \to Y$称为光滑的，如果对每一个点 $p \in X$，都存在一个包含 $p$ 的邻域$U \subset M$以及一个光滑函数$g: U \to N$使得在交集部分有$g|_{U \cap X} = f|_{U \cap X}$。换句话说，函数 $g$ 是函数 $f$ 的一个光滑延拓。若函数 $f$ 是双射、光滑的，并且它的逆映射也是光滑的，那么 $f$ 就称为一个微分同胚。
\subsection{局部描述}
在一定的温和条件下，可以通过局部方法来检验一个可微映射是否是微分同胚。这就是 Hadamard–Caccioppoli 定理[1]：

若$U$和$V$是$\mathbb{R}^n$中的两个连通开子集，且$V$是单连通的，那么一个可微映射$f: U \to V$如果满足以下条件，则$f$是一个微分同胚：$f$是适当映射；对于每个$x \in U$，其微分$Df_x: \mathbb{R}^n \to \mathbb{R}^n$是双射（因此是一个线性同构）。

一些说明：

对于函数$f$来说，集合$V$的单连通性是保证$f$在全局上可逆的关键条件（仅依赖于导数在每一点都是双射这一条件）。

例如，考虑复平方函数的“实数化”版本：
  $$
  \begin{cases}
  f:\mathbb{R}^2 \setminus \{(0,0)\} \;\to\; \mathbb{R}^2 \setminus \{(0,0)\}\\
  (x,y) \;\mapsto\; (x^2 - y^2,\, 2xy).
  \end{cases}~
  $$
在这种情况下，函数 $f$ 是满射，并且满足：
  $$
  \det Df_x = 4(x^2 + y^2) \neq 0.~
  $$
也就是说，虽然在每一点上微分$Df_x$都是双射，但函数$f$并不是可逆的，因为它不是单射，例如：$f(1,0) = (1,0) = f(-1,0)$。

由于在可微函数中，一个点处的微分
  $$
  Df_x: T_xU \;\to\; T_{f(x)}V~
  $$
是一个线性映射，因此当且仅当$Df_x$是双射时，它才有定义良好的逆映射$Df_x$的矩阵形式是一个$n \times n$的矩阵，其第$i$行、第$j$列的元素是：$\partial f_i/\partial x_j$。这个矩阵被称为雅可比矩阵，它常用于显式计算。

微分同胚必须发生在相同维度的流形之间。设映射$f$从维度 $n$ 的流形映射到维度 $k$ 的流形：若$n < k$，则$Df_x$永远不可能是满射；若$n > k$，则$Df_x$永远不可能是单射。因此，在这两种情况下，$Df_x$都无法成为双射。

如果$Df_x$在某点$x$上是双射，则称$f$在$x$附近是一个局部微分同胚，因为由连续性可知，$x$附近所有足够接近的点$y$上的$Df_y$也都是双射。

若从维度$n$到维度$k$的光滑映射$f$满足$Df$（或局部的 $Df_x$）是满射，则称 $f$ 是一个沉降（局部称作“局部沉降”）。若$Df$（或局部的$Df_x$）是单射，则称$f$是一个浸入（局部称作“局部浸入”）。

一个可微的双射函数不一定是微分同胚。例如：$f(x) = x^3$并不是 $\mathbb{R}$ 到 $\mathbb{R}$ 的微分同胚，因为它的导数在 $0$ 处为零，从而它的逆函数在 $0$ 处不可微。这是一个同胚但不是微分同胚的例子。

当$f$是可微流形之间的映射时，“微分同胚”比“同胚”是更强的条件：对于微分同胚，$f$及其逆映射都必须是可微的；对于同胚，$f$及其逆映射只需要连续即可。因此，每个微分同胚都是同胚，但并非每个同胚都是微分同胚。

一个映射$f: M \to N$是微分同胚，当且仅当它在局部坐标图下满足微分同胚的定义。更具体地：为$M$选取一组相容的坐标图，为$N$也选取一组相容的坐标图。令$\phi$ 是$M$上的坐标图，$\psi$是$N$上的坐标图，对应的像分别为$U$和$V$。若对所有满足$f(\phi^{-1}(U)) \subseteq \psi^{-1}(V)$的情况，组合映射$\psi f \phi^{-1}: U \to V$是一个微分同胚，则$f$是$M$到$N$的微分同胚。
\subsection{示例}
由于任何流形都可以在局部参数化，我们可以考虑一些从$\mathbb{R}^2 \to \mathbb{R}^2$的显式映射。
\begin{itemize}
\item 设：
$$
f(x, y) = \bigl(x^2 + y^3,\; x^2 - y^3 \bigr)~
$$
我们计算它的雅可比矩阵：
$$
J_f = 
\begin{pmatrix}
2x & 3y^2 \\
2x & -3y^2
\end{pmatrix}~
$$
该矩阵的行列式为零，当且仅当：$xy = 0$。这意味着，只有在远离$x$-轴和$y$-轴的区域，$f$才有可能是局部微分同胚。然而，$f$并不是双射，因为：$f(x, y) = f(-x, y)$，因此它不能是一个微分同胚。
\item 设：
$$
g(x, y) =
\bigl(
a_{0} + a_{1,0}x + a_{0,1}y + \cdots ,\;
b_{0} + b_{1,0}x + b_{0,1}y + \cdots
\bigr)~
$$
其中$a_{i,j}$和$b_{i,j}$是任意实数，省略的项都是关于$x$和$y$至少二次的高阶项。在$(0,0)$处的雅可比矩阵为：
$$
J_{g}(0,0) =
\begin{pmatrix}
a_{1,0} & a_{0,1} \\
b_{1,0} & b_{0,1}
\end{pmatrix}.~
$$
可以看出，当且仅当：
$$
a_{1,0} b_{0,1} - a_{0,1} b_{1,0} \neq 0~
$$
时，$g$在$(0,0)$处是一个局部微分同胚。换句话说，$g$的分量函数的线性项作为多项式必须线性无关。
\item 设：
$$
h(x, y) = \bigl(\sin(x^2 + y^2),\; \cos(x^2 + y^2)\bigr)~
$$
计算其雅可比矩阵：
$$
J_h =
\begin{pmatrix}
2x\cos(x^2 + y^2) & 2y\cos(x^2 + y^2)\\
-2x\sin(x^2 + y^2) & -2y\sin(x^2 + y^2)
\end{pmatrix}.~
$$
可以看到，这个雅可比矩阵的行列式在所有点都为零！事实上，$h$的像仅为单位圆：$\sin^2(x^2 + y^2) + \cos^2(x^2 + y^2) = 1$.
\end{itemize}
\subsubsection{表面形变}
在力学中，由应力引起的变换称为形变，这种变换可以用一个微分同胚来描述。设$f: U \to V$是两个表面$U$和$V$之间的一个微分同胚，那么其雅可比矩阵$Df$是一个可逆矩阵。事实上，要求对$U$中的每一点 $p$，其某个邻域内的雅可比矩阵$Df$都保持非奇异。假设在某个表面的坐标图下：$f(x, y) = (u, v)$.

全微分$u$的全微分为：
$$
du = \frac{\partial u}{\partial x} dx + \frac{\partial u}{\partial y} dy~
$$
同理，$v$的全微分也有类似表达式。

于是：$(du, dv) = (dx, dy) Df$,这表示一个线性变换，它以原点为固定点，可以看作是某种类型复数的作用。当$(dx, dy)$也被解释为同一类型的复数时，这个线性作用就对应于该复数平面中的复数乘法。由于复数乘法保留某种类型的角度（欧几里得角、双曲角或斜率角），且$Df$在表面上始终可逆，因此该类型的复数在整个表面上是一致的。由此，一个表面形变或两个表面之间的微分同胚具有一种保角性，即保留某种意义下的角度。
**微分同胚群（Diffeomorphism group）**

设 $M$ 是一个**可微流形**，并且它是**第二可数**且\*\*豪斯多夫（Hausdorff）\*\*的。

$M$ 的**微分同胚群**是指所有从 $M$ 到 $M$ 的 **$C^r$ 微分同胚**所组成的群，记作：

$$
\text{Diff}^r(M)
$$

当 $r$ 已知或不加特别说明时，通常简写为：

$$
\text{Diff}(M)。
$$

这是一个“庞大”的群，因为只要 $M$ **不是零维流形**，这个群就**不是局部紧致的（not locally compact）**。
